\documentclass[11pt,a4paper]{article}

\usepackage{amsmath,amssymb,amsthm}
\usepackage{hyperref}
\usepackage{microtype}
\usepackage{doi}

% ---------- theorem environments ----------
\newtheorem{theorem}{Theorem}[section]
\newtheorem{lemma}[theorem]{Lemma}
\newtheorem{corollary}[theorem]{Corollary}
\newtheorem{proposition}[theorem]{Proposition}
\newtheorem{definition}[theorem]{Definition}
\newtheorem{remark}[theorem]{Remark}
\newtheorem{problem}[theorem]{Problem}

% ---------- notation shortcuts ----------
\newcommand{\bstar}{\beta^{*}}
\newcommand{\dpair}{\delta_{\mathrm{pair}}}
\newcommand{\Heis}{\mathrm{Heis}_{3}(\mathbb{Z}/q\mathbb{Z})}
\newcommand{\Proj}{\Pi}
\newcommand{\chifield}{\chi}
\newcommand{\Obs}{\mathcal{O}}
\newcommand{\Ceff}{\mathcal{C}_{\mathrm{eff}}}
\newcommand{\rhoc}{\rho_{c}}
\newcommand{\rhoqc}{\rho_{q-c}}
\newcommand{\cchi}{c_{\mathrm{BI}}}
\newcommand{\Cay}{\mathrm{Cay}}


\title{%
  Equivariant Born--Infeld Parity and the Fibre-Identification Problem\\
  of the Non-Injective Projection\\[4pt]
  \large O18 --- Spectral Admissibility Sub-programme%
}
\author{J.~Beau\\[2pt]
\small Independent Researcher}
\date{2026}

\begin{document}
  \maketitle

  \begin{abstract}
    The identification of conjugate Weil blocks $\{c,\, q-c\}$ as the fibres of the
    non-injective projection $\Proj:\chifield\to\Obs$ was stated as a structurally motivated
    hypothesis in O16~\cite{Beau2026a20} and used in O17~\cite{Beau2026a21} to derive the
    pair-level observable and the exponent doubling $\dpair = 2\,\delta_{c} \approx 7.44$.
    The present paper determines exactly which part of that identification the Born--Infeld
    structure settles.
    We prove that the Born--Infeld action $S[\chifield]$ is even in $\chifield$, and that
    parity acts \emph{equivariantly} on the family of its response functionals: under the
    transported-probe comparison $(\chifield,\eta)\mapsto(-\chifield,-\eta)$, a response of
    derivative order $k$ transforms with the character $(-1)^{k}$, so even-order responses
    (the action itself among them) are preserved while odd-order responses --- including the
    Born--Infeld constitutive response --- reverse sign and generically separate
    $\chifield$ from $-\chifield$.
    We also prove, by an explicit countermodel, that evenness does not force the stronger
    fixed-probe (pointwise) indiscernibility of $\chifield$ and $-\chifield$: an even
    response can separate $\chifield$ from $-\chifield$ under one and the same perturbation.
    Consequently, the parity orbit $\{\chifield,-\chifield\}$ is a candidate fibre only
    relative to a supplied restriction of the observable family to even-order responses, and
    is not forced by the action alone under either comparison.
    We then state, as a typed classification problem, the data that a genuine fibre
    derivation must supply: the configuration space with its boundary conditions, the group
    of candidate transformations, the prior quotients, the separating family of responses,
    and a factorisation theorem connecting response equivalence to the fibres of $\Proj$.
    The absence-of-further-symmetries hypothesis (H-rank) is one clause of that problem, not
    a theorem and not yet a classifiable statement on its own.
    At the Weil level, the conjugation $\rhoqc = \overline{\rhoc}$ proves a spectral
    equivalence of the blocks $c$ and $q-c$; the identification of this equivalence with a
    physical projection fibre remains an open bridge.
    The pair-level observable of O16--O17 therefore keeps the status O16 gave it: a
    structurally motivated hypothesis, spectrally supported, awaiting a fibre theorem.
  \end{abstract}

  \textbf{Keywords.}
  fibre-level admissibility, Born--Infeld parity, parity covariance, parity involution,
  non-injective projection, Weil representation, Heisenberg Cayley graphs, pair observable,
  classification problem

  \clearpage
  \tableofcontents

  \clearpage

  \section{Context and scope}
    \label{sec:context}

    The Cosmochrony framework~\cite{Cosmochrony} describes physics as emerging from a
    pre-geometric substrate field $\chifield$ projected onto effective observables via a map
    $\Proj:\chifield\to\Obs$.
    A central structural motivation of the framework, developed in the companion
    paper~\cite{Beau2026c}, is that $\Proj$ is non-injective: injectivity would imply trivial
    holonomy and would therefore preclude any emergent gauge structure.
    Finite projective capacity is the capacity axiom [A-cap] of~\cite{Beau2026c}, motivated by
    non-injective projection with finite local distinguishability, and the Born--Infeld
    completion of the effective action is the saturation candidate compatible with this
    axiom~\cite{Beau2026c}.

    The O-series spectral admissibility sub-programme has investigated the capacity exponent
    $\delta$ governing projective dynamics on Heisenberg Cayley graphs $\Cay(G_q,S_q)$ with
    $G_q = \Heis$ for prime $q$.
    Papers O12--O15~\cite{Beau2026a16,Beau2026a17,Beau2026a18,Beau2026a19} established the
    exact block-level exponent $\hat{\delta}_{\mathrm{exact}} \approx 3.72$, structurally
    below the phenomenological target $\delta \in [7.4, 10.6]$.
    O16~\cite{Beau2026a20} identified the resolution: the physically relevant observable
    lives at the level of conjugate pairs $\{c, q-c\}$, yielding $\dpair = 2\,\delta_c
  \approx 7.44$.
    O17~\cite{Beau2026a21} provided the structural derivation showing that the Gram--Schmidt
    dynamics of blocks $c$ and $q-c$ are identical, establishing $\delta_{q-c} = \delta_c$
    and that the factor $r(c,q)$ in O16 is a normalisation artefact, not an invariant.

    O17 Section~5.2 identifies as the primary open problem of the fibre-level programme the
    derivation of the identification $c \leftrightarrow q-c$ as a necessary consequence of
    the projection structure $\Proj$ and the Born--Infeld constraint, rather than a
    structurally motivated hypothesis.
    The present paper examines that problem and establishes exactly which part of it the
    Born--Infeld structure settles.

    \paragraph{Scope and what this paper does and does not establish.}
      O18 establishes four things:
      (i) the Born--Infeld action is even in $\chifield$;
      (ii) parity acts equivariantly on the Born--Infeld response family: responses of
      derivative order $k$ transform with the character $(-1)^{k}$ under the
      transported-probe comparison, so even-order responses are preserved and odd-order
      responses reverse sign;
      (iii) evenness does \emph{not} force fixed-probe indiscernibility of $\chifield$ and
      $-\chifield$: an explicit even countermodel separates them under a common
      perturbation, so the parity orbit is not forced into the fibres by the action alone;
      (iv) the derivation of the fibre structure of $\Proj$ is a classification problem
      whose required data --- transformation group, prior quotients, boundary conditions,
      response family (including its restriction, if any, to even-order responses), and a
      factorisation theorem for $\Proj$ --- are stated here explicitly and remain
      unsupplied.
      At the Weil level, the correspondence $c \leftrightarrow q-c$ is a structural
      candidate supported by the spectral equivalence $\rhoqc=\overline{\rhoc}$ of
      O17~\cite{Beau2026a21}; its identification with a physical projection fibre is an
      open bridge.
      O18 does \emph{not} derive the fibre structure of $\Proj$ from the Born--Infeld
      action, and no downstream result may cite it for such a derivation.

  \section{Abstract level: parity from Born--Infeld}
  \label{sec:abstract}

  \subsection{Imported results on non-injectivity}
    \label{ssec:noninjectivity}

    We recall two structural inputs from the companion paper~\cite{Beau2026c} that are used
    here without re-derivation.

    \begin{enumerate}
      \item \textbf{Non-injectivity of $\Proj$}~\cite{Beau2026c}.
      If $\Proj$ were injective and compatible with relational locality, the induced holonomy
      would be trivial, which would preclude emergent gauge structure.
      Since electromagnetism must arise as an effective U(1) fibre symmetry, this argument
      motivates taking $\Proj$ to be non-injective.
      Non-trivial fibres $\Proj^{-1}(y)$ are therefore structurally motivated rather than an
      arbitrary auxiliary assumption.

      \item \textbf{Capacity axiom and Born--Infeld candidate}~\cite{Beau2026c}.
      Finite projective transport capacity is the capacity axiom [A-cap] of~\cite{Beau2026c},
      an explicit hypothesis motivated by non-injective projection with finite local
      distinguishability.
      The Born--Infeld effective action is the saturation candidate compatible with this
      axiom, analytic in the weak-field limit and saturating the constitutive relation; it is
      the form adopted throughout this paper.
      The effective action takes the form
      \begin{equation}
        S_{\mathrm{BI}}[\chifield]
        = \beta^{2} \int \!\mathrm{d}^{4}x
        \left(1 - \sqrt{1 - \frac{F^{2}}{\beta^{2}}}\right),
        \qquad F = D\chifield,
        \label{eq:SBI}
      \end{equation}
      where $\beta$ is the saturation scale and $F^{2} = F_{\mu\nu}F^{\mu\nu}$.
    \end{enumerate}

  \subsection{Evenness of the Born--Infeld action}
    \label{ssec:parity}

    \begin{lemma}[Parity of $S_{\mathrm{BI}}$]
      \label{lem:parity}
      $S_{\mathrm{BI}}[-\chifield] = S_{\mathrm{BI}}[\chifield]$.
    \end{lemma}

    \begin{proof}
      Under $\chifield \mapsto -\chifield$, the field gradient transforms as
      $F = D\chifield \mapsto -D\chifield = -F$.
      Therefore $F^{2} = F_{\mu\nu}F^{\mu\nu} \mapsto (-F)_{\mu\nu}(-F)^{\mu\nu} = F^{2}$.
      Since the integrand in~\eqref{eq:SBI} depends on $\chifield$ only through $F^{2}$,
      the action is unchanged.
    \end{proof}

    \begin{remark}
      Lemma~\ref{lem:parity} uses only the quadratic dependence on $F$ in the argument of the
      square root, which is the defining property of the Born--Infeld structure.
      It does not depend on any discrete group structure and holds for all smooth $\chifield$.
    \end{remark}

  \subsection{Two notions of indiscernibility}
    \label{ssec:indiscern}

    A fibre statement compares the responses of two configurations to perturbations.
    There are two inequivalent ways to organise that comparison, and the distinction carries
    the main result of this section.

    \begin{definition}[Fixed-probe BI-indiscernibility]
      \label{def:indiscern}
      Two configurations $\chifield_{1}, \chifield_{2}$ are \emph{fixed-probe
      BI-indiscernible} if, for every BI-admissible perturbation $\eta$ (i.e.\ $\|D\eta\|$
      bounded and $\|D(\chifield_{i} + \epsilon\eta)\| \leq \cchi$ for small
      $\epsilon > 0$) and every effective response functional $R$ constructed solely from
      $S_{\mathrm{BI}}$ and its derivatives,
      \begin{equation}
        R[\chifield_{1} + \epsilon\eta]
        = R[\chifield_{2} + \epsilon\eta]
        \quad\text{for all admissible } \eta \text{ and all sufficiently small } \epsilon.
        \label{eq:indiscern}
      \end{equation}
      The same perturbation $\eta$ is applied to both configurations.
    \end{definition}

    \begin{proposition}[Evenness does not force fixed-probe indiscernibility]
      \label{prop:no-pointwise}
      There exist even response functionals $R$, configurations $\chifield$, and admissible
      perturbations $\eta$ with
      $R[\chifield + \epsilon\eta] \neq R[-\chifield + \epsilon\eta]$ for all sufficiently
      small $\epsilon > 0$.
      Consequently, $\chifield$ and $-\chifield$ are in general \emph{not} fixed-probe
      BI-indiscernible, and evenness of $S_{\mathrm{BI}}$ does not place the parity orbit
      $\{\chifield,-\chifield\}$ inside the fibres of any projection defined by
      Definition~\ref{def:indiscern}.
    \end{proposition}

    \begin{proof}
      Take a one-dimensional toy model with configurations $x \in \mathbb{R}$ and the even
      response $R(x) = x^{2}$, which is the toy action itself.
      For $\chifield = 1$, $\eta = 1$, and any sufficiently small $\epsilon > 0$,
      \[
        R(1+\epsilon) = (1+\epsilon)^{2}
        \neq (1-\epsilon)^{2} = R(-1+\epsilon).
      \]
      The probe $\eta = 1$ is admissible in the sense of Definition~\ref{def:indiscern}
      whenever $\cchi \geq 1 + \epsilon$, which holds for all sufficiently small $\epsilon$
      once $\cchi > 1$.
      For an even response, evenness gives only the reindexed family equality
      $R[-\chifield + \epsilon\eta] = R[\chifield - \epsilon\eta]$; equality of a family
      after the bijective reindexing $\eta \mapsto -\eta$ does not give equality at the same
      index $\eta$, and the countermodel shows that it fails.
      For the field version, choose an admissible perturbation with
      $D\chifield \cdot D\eta \neq 0$: the cross term in
      $(D\chifield + \epsilon D\eta)^{2}$ changes sign between $\chifield$ and $-\chifield$,
      so the action itself --- a response functional built solely from $S_{\mathrm{BI}}$ ---
      separates the two configurations under the common probe $\eta$.
    \end{proof}

    \begin{definition}[Transported-probe comparison; derivative order of a response]
      \label{def:equivariant}
      Let $g$ be a transformation acting linearly on configurations and perturbations.
      The \emph{transported-probe comparison} of $\chifield$ with $g\chifield$ evaluates
      each response at $(\chifield + \epsilon\eta)$ and at
      $(g\chifield + \epsilon\, g\eta)$: the probe is transported by the same
      transformation as the configuration.
      A response functional \emph{of derivative order $k$} is a $k$-th functional
      derivative of $S_{\mathrm{BI}}$ evaluated at the configuration and contracted with
      $k$ fixed admissible test directions; the test directions are part of the response's
      definition and are \emph{not} transported by $g$ (only the configuration and the
      probe are).
      The responses of Definition~\ref{def:indiscern} are generated by these homogeneous
      responses, with $k = 0$ giving the action-level responses; a general response is a
      polynomial in homogeneous ones, and since the character below is multiplicative, a
      product of responses of orders $k_1,\dots,k_m$ transforms with $(-1)^{k_1+\dots+k_m}$
      and a sum inherits the behaviour of its terms.
    \end{definition}

    \begin{proposition}[Parity covariance of the response family]
      \label{prop:equivariance}
      Let $R$ be a response functional of derivative order $k$.
      Under the parity map $g:\chifield\mapsto-\chifield$, $\eta\mapsto-\eta$,
      \begin{equation}
        R[-\chifield - \epsilon\eta] = (-1)^{k}\, R[\chifield + \epsilon\eta]
        \qquad\text{for all admissible } \eta.
        \label{eq:parity-equivariance}
      \end{equation}
      Parity therefore acts equivariantly on the response family through the character
      $(-1)^{k}$: even-order responses --- the action itself among them --- take equal
      values at $(\chifield,\eta)$ and $(-\chifield,-\eta)$, while odd-order responses
      reverse sign and generically separate $\chifield$ from $-\chifield$ under the
      transported probe.
    \end{proposition}

    \begin{proof}
      $-\chifield - \epsilon\eta = -(\chifield + \epsilon\eta)$, so it suffices to show
      $R[-\psi] = (-1)^{k} R[\psi]$ for the $k$-th variation of $S_{\mathrm{BI}}$.
      Differentiating the evenness identity $S_{\mathrm{BI}}[-\psi] = S_{\mathrm{BI}}[\psi]$
      of Lemma~\ref{lem:parity} $k$ times along admissible directions produces one sign per
      differentiation of the inner argument, hence the factor $(-1)^{k}$.
      Admissibility is preserved because $\|D(-\psi)\| = \|D\psi\|$.
      The first-order instance is the Born--Infeld constitutive response, which is odd:
      $-\chifield$ is the opposite-sign configuration, and any response resolving a
      non-vanishing first variation distinguishes it from $\chifield$ (separation fails
      only where the contracted first variation vanishes).
    \end{proof}

    \begin{remark}[What covariance does and does not give]
      \label{rem:equivariance-scope}
      Proposition~\ref{prop:equivariance} exhibits parity as a symmetry of the Born--Infeld
      \emph{action}, acting equivariantly on responses with the character $(-1)^{k}$.
      Indiscernibility of $\chifield$ and $-\chifield$ under the transported-probe
      comparison therefore holds on the even-order part of the response family, and fails
      generically on the odd-order part (wherever the odd-order response does not vanish).
      The parity orbit $\{\chifield,-\chifield\}$ is thus a candidate fibre only relative
      to a \emph{supplied} restriction of the observable family to even-order responses;
      that restriction is part of the response-family datum (D3) of
      Problem~\ref{prob:classification} and is not derived here.
      Given such a restriction together with a supplied group action containing parity, the
      orbits define a quotient of configuration space, and a projection may consistently be
      \emph{chosen} to factor through that quotient.
      Covariance does not prove that the physical projection $\Proj$ has these orbits as
      fibres: that requires a factorisation theorem connecting response equivalence to
      $\Proj$, which is part of the classification problem of
      Section~\ref{ssec:classification} and is not available.
      By Proposition~\ref{prop:no-pointwise}, the fixed-probe route to the same conclusion
      is closed, so no current result forces the parity orbit into the fibres of $\Proj$.
    \end{remark}

  \subsection{The fibre-identification problem, typed}
    \label{ssec:classification}

    A derivation of the fibre structure of $\Proj$ from the effective action must solve a
    classification problem whose data we now state explicitly.

    \begin{problem}[Fibre identification for the Born--Infeld projection]
      \label{prob:classification}
      Supply the following data and prove the following theorem.
      \begin{enumerate}
        \item[(D1)] The configuration space $\mathcal{C}_{\chi}$, with its topology and
        boundary conditions.
        \item[(D2)] The group $\mathcal{G}$ of candidate transformations acting on
        $\mathcal{C}_{\chi}$ and on probes, together with the prior quotients already taken
        (spacetime, gauge, internal, field redefinitions).
        \item[(D3)] The complete separating family $\mathcal{R}$ of effective response
        functionals, with the notion of comparison (fixed-probe or transported-probe) fixed
        once for the whole family, and with any restriction to even-order responses stated
        explicitly.
        \item[(D4)] The factorisation theorem: the equivalence relation on
        $\mathcal{C}_{\chi}$ induced by $\mathcal{R}$-indiscernibility coincides with the
        fibre partition of the physical projection $\Proj$.
      \end{enumerate}
      Given (D1)--(D4), classify the effective symmetries of $S_{\mathrm{BI}}$ within
      $\mathcal{G}$ and determine the fibre of $\Proj$ through each configuration.
    \end{problem}

    \begin{remark}[Status of the absence-of-further-symmetries hypothesis]
      \label{rem:hrank-status}
      The hypothesis that parity is the only effective symmetry of $S_{\mathrm{BI}}$
      compatible with BI-admissible observables --- hypothesis (H-rank) in downstream
      papers --- is a candidate answer to the classification clause of
      Problem~\ref{prob:classification}.
      Before (D1)--(D3) are supplied, it is not a classifiable mathematical statement: the
      class of transformations over which ``only'' quantifies is not defined.
      It is therefore not an irreducible modelling choice either; it is a clause of a
      problem not yet posed in full.
      Two structural facts show that the clause has genuine content.
      First, the adopted action depends on $\chifield$ only through $D\chifield$, so the
      constant shifts $\chifield \mapsto \chifield + a$ are symmetries of
      $S_{\mathrm{BI}}$ unless boundary conditions or a prior quotient in (D1)--(D2) remove
      them; parity is thus not automatically the only candidate.
      Second, an equivalence relation can have non-trivial classes without being the orbit
      relation of any globally defined group action, so response indiscernibility of a pair
      of configurations does not by itself produce a global symmetry.
    \end{remark}

  \section{Concrete level: Weil realisation in $\Heis$}
  \label{sec:concrete}

  \subsection{The setting}
    \label{ssec:setting}

    We work in the discrete Heisenberg group $G_q = \Heis$ for prime $q$, with the Cayley
    graph $\Cay(G_q, S_q)$ where $S_q = \{X^{\pm 1}, Y^{\pm 1}\}$ and $X, Y$ are the
    standard generators.
    The Weil representations $\rhoc$ of $G_q$ are parametrised by the character index
    $c \in (\mathbb{Z}/q\mathbb{Z})^{\times}$.
    The Gram--Schmidt redundancy $\sigma_c(n)$ at BFS shell $n$ is the key spectral
    observable established in O12--O15~\cite{Beau2026a16,Beau2026a17,Beau2026a18,Beau2026a19}.

  \subsection{Imported result from O17}
    \label{ssec:O17}

    The following result is Theorem~2.1 and Corollary~2.3 of O17~\cite{Beau2026a21}.

    \begin{theorem}[O17 Theorem 2.1 and Corollary 2.3]
      \label{thm:O17}
      The conjugation identity $\rhoqc = \overline{\rhoc}$ together with the norm-invariance of
      the Gram--Schmidt process implies that the raw Gram--Schmidt redundancy counts are
      exactly identical for $\rhoc$ and $\rhoqc$ when the initial supports coincide, and
      structurally isomorphic in general.
      In particular, $\delta_{q-c} = \delta_c$ for all admissible $c$.
    \end{theorem}

  \subsection{The conjugate pair as a candidate fibre: an open bridge}
    \label{ssec:identification}

    The conjugation $\rhoqc = \overline{\rhoc}$ makes the pair $\{c, q-c\}$ the natural
    candidate realisation of the abstract parity orbit in the Weil setting: complex
    conjugation reverses representation phases, which is how a sign inversion of an
    effective amplitude would manifest spectrally, and Theorem~\ref{thm:O17} proves that the
    two blocks are dynamically indistinguishable under Gram--Schmidt.

    What Theorem~\ref{thm:O17} proves is a \emph{spectral equivalence} of the blocks $c$
    and $q-c$.
    The further statement that the two blocks form one fibre of the physical projection
    $\Proj$ --- i.e.\ the correspondence $\chifield \leftrightarrow c$,
    $-\chifield \leftrightarrow q-c$ --- is an identification between a
    representation-theoretic label and a substrate-field sign.
    No derivation of this correspondence from $\chifield$-space variables is available, and
    complex conjugation alone does not supply it: equality of Gram--Schmidt norms for
    conjugate representations proves that the two blocks are spectrally equivalent, not
    that $\Proj$ maps them to the same observable.

    \begin{remark}[Status of the Weil-level identification]
      \label{rem:identification}
      The correspondence $\chifield \leftrightarrow c$, $-\chifield \leftrightarrow q-c$ is
      an open bridge: a structural candidate, spectrally supported by
      Theorem~\ref{thm:O17}, whose derivation would require solving
      Problem~\ref{prob:classification} together with an explicit construction relating
      $\chifield$-space parity to the Weil character index.
      Any downstream use of the pair $\{c,q-c\}$ as a fibre of $\Proj$ must carry this
      status explicitly.
    \end{remark}

  \section{Implications and relation to O16--O17}
  \label{sec:implications}

  \subsection{Status of the pair-level programme}
    \label{ssec:requalification}

    O16~\cite{Beau2026a20} introduced the pair observable $\sigma_{\mathrm{pair}}(n) =
  \sigma_c(n) \cdot \sigma_{q-c}(n)$ and stated the identification of $\{c, q-c\}$ with
    fibres of $\Proj$ as a structurally motivated hypothesis (O16 Remark~6.2).
    O17~\cite{Beau2026a21} derived the equality $\delta_{q-c} = \delta_c$ from
    $\rhoqc = \overline{\rhoc}$ and identified $r(c,q)$ as a normalisation artefact.
    These spectral results stand.
    The fibre identification itself keeps the status O16 gave it: a structurally motivated
    hypothesis.
    The present paper adds the covariance result (Proposition~\ref{prop:equivariance}),
    which makes the parity orbit the distinguished candidate relative to the even-order
    part of the response family, and the no-go of
    Proposition~\ref{prop:no-pointwise}, which shows that the fixed-probe route cannot
    upgrade the hypothesis to a theorem.
    The upgrade, if available, must come from a solution of
    Problem~\ref{prob:classification}.

    \begin{remark}[Complementarity with O15--O17]
O15~\cite{Beau2026a19} established that no reweighting of block-level data can recover
the fibre-level exponent.
O16--O17 established the correct observable and exponent at the pair level, conditional
on the fibre hypothesis.
O18 (the present paper) delimits what the Born--Infeld structure itself contributes to
that hypothesis: parity covariance of the response family, and no more.
    \end{remark}

  \subsection{Effect on $\delta_{\mathrm{pair}} \to \bstar$}
    \label{ssec:dbetastar}

    The structural relation $\bstar = 1/(\delta + 1/2)$, used throughout the O-series, is
    evaluated at $\delta = \dpair \approx 7.44$ in O17 Section~4 on the basis of the fibre
    identification.
    That evaluation is therefore conditional on the same hypothesis as the identification
    itself, and it inherits the open-bridge status of
    Remark~\ref{rem:identification}.
    The aggregation no-go of O15 Corollary~4.8~\cite{Beau2026a19} stands: no block-level
    reweighting recovers the fibre-level exponent, so the pair level remains the only
    candidate carrier of the phenomenological exponent, conditional on the fibre
    hypothesis.

  \subsection{Open directions}
    \label{ssec:open}

    The present paper leaves four directions open.

    \textbf{The classification problem.}
    Problem~\ref{prob:classification} is the central open item: supply (D1)--(D4) and
    classify the effective symmetries.
    Until it is solved, the fibre structure of $\Proj$ is a hypothesis in every downstream
    use.

    \textbf{The Weil-level bridge.}
    A derivation of the correspondence $\chifield \leftrightarrow c$,
    $-\chifield \leftrightarrow q-c$ from $\chifield$-space variables
    (Remark~\ref{rem:identification}).

    \textbf{Canonical normalisation.}
    The pipeline normalisation factor $D(c, b_1, b_2, n)$ is currently implicit in the
    O12/O13 implementation.
    Defining a canonical normalisation independent of $(b_1, b_2)$ would eliminate $r$
    entirely from the pair observable, producing a version of $\sigma_{\mathrm{pair}}$ that
    is genuinely independent of the measurement convention.
    This is a technical, not conceptual, task and is the subject of O19.

    \textbf{Upper bound of the phenomenological window.}
    Numerical exploration (post-O16) shows that the spectrum of $\dpair$ has a hard lower
    bound near $7.4$ and extends above $10.6$.
    The mechanism selecting $[7.4, 10.6]$ as the physical sub-spectrum is not yet identified.
    It appears to be a dynamical persistence criterion rather than a spectral one, and is the
    subject of O20.

  \section{Summary}
  \label{sec:summary}

  \begin{enumerate}
    \item The Born--Infeld action $S_{\mathrm{BI}}[\chifield]$ is even in $\chifield$
    (Lemma~\ref{lem:parity}).
    This follows directly from the dependence of $S_{\mathrm{BI}}$ on $F^2 = (D\chifield)^2$.

    \item Parity acts equivariantly on the Born--Infeld response family with the
    character $(-1)^{k}$: even-order responses are preserved under the transported-probe
    comparison and odd-order responses reverse sign
    (Proposition~\ref{prop:equivariance}).
    This makes the orbit $\{\chifield, -\chifield\}$ the distinguished candidate fibre
    relative to the even-order part of the family.

    \item Evenness does not force fixed-probe indiscernibility: the even response
    $R(x) = x^{2}$ separates $\chifield = 1$ from $-\chifield = -1$ under the common probe
    $\eta = 1$ (Proposition~\ref{prop:no-pointwise}).
    The parity orbit is therefore supplied by a comparison convention together with an
    even-order restriction of the observables, not forced by the action.

    \item The derivation of the fibre structure of $\Proj$ is the typed classification
    problem stated in Problem~\ref{prob:classification}; its data (D1)--(D4) are not
    supplied by the current corpus.
    The absence-of-further-symmetries hypothesis (H-rank) is one clause of that problem
    (Remark~\ref{rem:hrank-status}).

    \item In the Weil realisation of $\Heis$, the conjugate pair $\{c, q-c\}$ is the
    candidate realisation of the parity orbit, spectrally supported by
    $\rhoqc = \overline{\rhoc}$ and $\delta_{q-c} = \delta_c$
    (Theorem~\ref{thm:O17}); the identification with a fibre of $\Proj$ is an open bridge
    (Remark~\ref{rem:identification}).
    The pair observable of O16--O17 keeps its status of structurally motivated hypothesis.
  \end{enumerate}

  \bibliographystyle{plain}
      \bibliography{cosmochrony-bibliography}

\end{document}
